\documentclass{article}
\usepackage[utf8]{inputenc}
\usepackage{booktabs}
\title{Lightweight Building Blocks for Efficient Sequence Classification}
\author{A. Example}
\begin{document}
\maketitle

\begin{abstract}
We study which standard deep learning components matter most for small text classifiers.
Our compact model improves accuracy by 3.2 points over a strong baseline on the ten-class news benchmark.
It also trains twice as fast as the baseline.
\end{abstract}

\section{Introduction}
Self-attention has replaced recurrence in many sequence models.
The Transformer reached 28.4 BLEU on the WMT 2014 English-to-German translation task \cite{vaswani2017attention}.
It was trained for 3.5 days on eight P100 GPUs \cite{vaswani2017attention}.
Residual connections make very deep networks trainable, and residual networks with 152 layers won first place in the ILSVRC 2015 classification task \cite{he2016deep}.
Pre-trained bidirectional encoders outperform human annotators on every GLUE task \cite{devlin2019bert}.

\section{Background}
Adam is an optimizer that requires a large amount of memory compared with stochastic gradient descent \cite{kingma2015adam}.
Batch normalization reaches the same accuracy as the original network with 14 times fewer training steps \cite{ioffe2015batch}.
Deep residual networks reached a 3.57\% top-5 error on the ImageNet test set \cite{he2016deep}, while the Transformer reported 41.8 BLEU on English-to-French \cite{vaswani2017attention}.
The BERT large model has 340 million parameters \cite{devlin2019bert}.

\section{Method}
Our model combines a two-layer self-attention encoder with residual connections and batch normalization.
We optimise it with Adam using a learning rate of 0.001 and the default moment estimates \cite{kingma2015adam}.

\section{Results}
Table~\ref{tab:main} summarises the results on the ten-class news benchmark.
Our compact model reaches 91.2\% accuracy, while the baseline reaches 88.0\%.
Training our model takes 2.1 hours compared with 4.0 hours for the baseline.
On the out-of-domain test set, accuracy drops to 79.5\% for our model and 78.9\% for the baseline.

\begin{table}[h]
\centering
\caption{Accuracy on the ten-class news benchmark.}
\label{tab:main}
\begin{tabular}{lcc}
\toprule
Model & Accuracy (\%) & Training time (h) \\
\midrule
Baseline & 88.0 & 4.0 \\
Ours & 91.2 & 2.1 \\
\bottomrule
\end{tabular}
\end{table}

\section{Conclusions}
Our compact model improves accuracy by more than 5 points over the baseline.
The gains hold equally well on out-of-domain data and generalise to every language.
Training is roughly twice as fast as the baseline, which makes the model practical on modest hardware.

\bibliographystyle{plain}
\bibliography{refs}
\end{document}
